% tested with pdflatex, tectonic emits warnings
% problem: form values don't appear in print
\documentclass[11pt,a4paper]{article}
\usepackage[utf8]{inputenc}
\usepackage[T1]{fontenc}
\usepackage{hyperref}
\setlength{\parindent}{0pt}
% form fields use sans serif font, adapt document font to match
\usepackage[scaled]{helvet}
\renewcommand{\familydefault}{\sfdefault}
\hypersetup{pdfborder={0 0 0},colorlinks=false}
\def\DefaultOptionsofText{
bordercolor={0 0 0},
borderwidth=0.0pt,
backgroundcolor={1 1 1},
}
%\def\DefaultOptionsofCheckBox{bordercolor={0 0 0}, borderwidth=0.4pt, backgroundcolor={1 1 1}}
\def\DefaultOptionsofCheckBox{bordercolor={0 0 0}}
\usepackage{microtype}
\usepackage[ngerman]{babel}
\usepackage[margin=2.5cm]{geometry}
\renewcommand\thesubsubsection{\arabic{subsubsection}}

\begin{document}

\section*{Informationsschreiben und Einwilligungserklärung zur Masterarbeit}

\begin{center}
\TextField[name=titel,width=17cm,charsize=17pt,multiline=true,default={Titel Ihrer Arbeit}]{}
\end{center}

\section*{Projektbeteiligte und Kontaktpersonen}

\subsubsection*{Durchführende Person}
\TextField[name=name,width=17cm,default={Vorname Name}]{}\\
Studiengang: \TextField[name=studiengang,width=6cm,default={Medizininformatik M. Sc.}]{} \\
\TextField[name=email,width=17cm,default={...@studserv.uni-leipzig.de}]{} \\

\subsubsection*{Betreuung}
\TextField[name=betreuer,width=17cm,default={Titel Vorname Nachname Betreuer:in}]{}\\
\TextField[name=inst,width=17cm,default={Institut für Medizinische Informatik, Statistik und Epidemiologie}]{}
\TextField[name=uni,width=17cm,default={Universität Leipzig}]{}
\TextField[name=emailbet,width=17cm,default={Email-Adresse Betreuer:in}]{} \\

\vspace{2em}

Liebe/r Teilnehmer/in,

dieses Schreiben dient dazu, Sie über den Zweck, den Ablauf und den Umgang mit Ihren Daten, die durch die Interviews erhoben werden, zu informieren.
Bitte lesen Sie den nachfolgenden Text aufmerksam durch.
Sollten Sie mit der Teilnahme an einem Interview einverstanden sein, unterschreiben Sie gerne die Einverständniserklärung und schicken Sie diese per Mail an die oben genannten Kontaktdaten der durchführenden Person.
Falls Sie noch weitere Fragen haben, kontaktieren Sie mich gerne jederzeit.

\section{Informationen über die Masterarbeit}
\TextField[name=info,width=17cm,multiline=true,
default={1-2 Sätze zur Zielsetzung Ihrer Arbeit und der Interviews.
Die Datenerhebung erfolgt im Rahmen eines halbstrukturierten, leitfadengestützten Interviews,
das nach Zustimmung der Teilnehmenden digital aufgezeichnet wird.
Die Audioaufnahme wird anschließend transkribiert.
Die Datenspeicherung erfolgt digital auf gesicherten, passwortgeschützten Datenträgern am Institut für Medizinische Informatik, Statistik und Epidemiologie.
}]{}

\section{Datenschutz}

\subsection*{Allgemeine Informationen}

\subsubsection{Verantwortliche Stelle gemäß Art. 4 Abs. 7 DSGVO}

Universität Leipzig\\
Medizinische Fakultät\\
Liebigstr. 27 A\\
04103 Leipzig \\
Telefon: +49 341 97-15930\\
E-Mail: \href{mailto:dekanat@medizin.uni-leipzig.de}{dekanat@medizin.uni-leipzig.de}

\subsubsection*{Kontaktdaten des Datenschutzbeauftragten}

Ronald Speer\\
Datenschutzbeauftragter\\
Medizinische Fakultät der Universität Leipzig\\
Philipp-Rosenthal-Str. 27\\
04103 Leipzig\\
Telefon: +49 341 97-16105\\
E-Mail: \href{mailto:dsbmf@medizin.uni-leipzig.de}{dsbmf@medizin.uni-leipzig.de}

\subsubsection{Sie haben folgende Rechte hinsichtlich Ihrer personenbezogenen Daten}

\begin{itemize}
\item Recht auf Auskunft,
\item Recht auf Berichtigung,
\item Recht auf Löschung,
\item Recht auf Einschränkung der Verarbeitung,
\item Recht auf Widerspruch gegen die Verarbeitung,
\item Recht auf Datenübertragbarkeit
\end{itemize}

sofern diese Rechte nicht durch Gesetze im Rahmen der Datenverarbeitung zu wissenschaftlichen Zwecken eingeschränkt sind.
Sie haben zudem das Recht, sich bei einer Datenschutz-Aufsichtsbehörde über die Verarbeitung Ihrer personenbezogenen Daten in unserem Unternehmen zu beschweren.

\subsubsection{Bei Fragen zur Datenverarbeitung können Sie sich an die durchführende Person oder die betreuenden Personen wenden.}

\subsubsection{Die Projektdurchführenden werden alle angemessenen Schritte unternehmen, um den Schutz Ihrer Daten gemäß Datenschutz-Grundverordnung (DS-GVO) und anderen Gesetzen zu gewährleisten.}
Die Daten sind gegen unbefugten Zugriff gesichert.
Die personenbezogenen Daten werden anonymisiert, sobald dies nach dem Forschungszweck möglich ist, es sei denn berechtigte Interessen der betroffenen Person stehen dem entgegen.
Bis dahin werden die Merkmale gesondert gespeichert, mit denen Einzelangaben über persönliche oder sachliche Verhältnisse einer bestimmten oder bestimmbaren Person oder Institution zugeordnet werden können.
Sie werden mit den Einzelangaben nur zusammengeführt, soweit der Forschungszweck dies erfordert.

\subsubsection{Die Verarbeitung ihrer personenbezogenen Daten erfolgt nur, sofern Sie ausdrücklich in die Erhebung und Verarbeitung Ihrer Daten eingewilligt haben.}

\subsubsection{Sie haben das Recht, Ihre Einwilligung in die Verarbeitung Ihrer personenbezogenen Daten jederzeit zu widerrufen.}
Durch den Widerruf der Einwilligung wird die Rechtmäßigkeit, der aufgrund der Einwilligung bis zum Widerruf erfolgten Verarbeitung nicht berührt.

\subsection*{Projektspezifische Informationen}
\setcounter{subsubsection}{0}

Im Folgenden informiere ich Sie über die Erhebung personenbezogener Daten im Zusammenhang mit der oben genannten Masterarbeit.

\subsubsection{Beschreibung und Umfang der Datenverarbeitung}

Zur Erfüllung der Zwecke der Masterarbeit werden personenbezogene Daten der Teilnehmer:innen erhoben.
Diese beinhalten insbesondere:

\begin{itemize}
\item \TextField[name=iden,width=15cm,default={Identifikationsdaten (z.B. Name, Kontaktdaten)}]{}
\item \TextField[name=beruf,width=15cm,default={Berufliche Angaben (z.B. Funktion, Tätigkeitsbereich)}]{}
\item \TextField[name=audio,width=15cm,default={Audioaufnahmen}]{}
\item \TextField[name=inhalt,width=15cm,default={Inhaltliche Daten aus dem Interview, insbesondere …}]{}
\end{itemize}

Der Zugriff auf diese Daten ist ausschließlich auf die durchführenden und betreuenden Personen beschränkt.
Personenbezogene sowie eindeutig zuordenbare institutionsbezogene Angaben werden nach der Erhebung pseudonymisiert.
Die Zuordnung zwischen den verwendeten Pseudonymen und den tatsächlichen Personen erfolgt ausschließlich über eine getrennt aufbewahrte Identifikationsliste.
Während dieses Zeitraums besteht daher eine theoretische Re-Identifikationsmöglichkeit.
Für Auswertung, Zitation und wissenschaftliche Darstellung werden ausschließlich pseudonymisierte bzw. nach Löschung der Identifikationsliste anonymisierte Aussagen verwendet.

\subsubsection{Rechtsgrundlage für die Datenverarbeitung}

Rechtsgrundlage ist die Einwilligungserklärung gem. Art. 6 Abs. 1 lit a) und Art. 9 Abs. 2 lit. a) DS-GVO.

\subsubsection{ Zweck der Datenverarbeitung}

Die Verarbeitung der erhobenen Daten erfolgt zur Bearbeitung der Masterarbeit.
Darüber hinaus können Inhalte und Aussagen für wissenschaftliche Veröffentlichungen genutzt werden, sofern diese keine Rückschlüsse auf einzelne Personen oder Institutionen zulassen.


\subsubsection{Dauer der Speicherung}

Die Audioaufnahmen werden nach der Transkription gelöscht.
Die pseudonymisierten Transkripte werden entsprechend den Empfehlungen der Deutschen Forschungsgemeinschaft (DFG) zur Sicherung guter wissenschaftlicher Praxis für einen Zeitraum von 15 Jahren gespeichert.
 Die Identifikationsliste wird bis zu diesem Zeitpunkt getrennt aufbewahrt.
Nach Ablauf der Speicherfrist werden die Identifikationsliste sowie sämtliche personenbezogenen Daten vollständig und unwiderruflich gelöscht.
Die verbleibenden Transkripte sind ab diesem Zeitpunkt anonymisiert und eine Zuordnung zu einzelnen Personen ist nicht mehr möglich.

\subsubsection{Widerspruchs- und Beseitigungsmöglichkeit}

Die Teilnahme ist freiwillig.
Sie können die Einwilligung jederzeit ohne Angabe von Gründen widerrufen, ohne dass Ihnen daraus Nachteile entstehen.
Der Widerruf kann per E-Mail an die im Schreiben oben genannte durchführende Person gerichtet werden.
Im Fall des Widerrufs werden Ihre Daten nicht weiterverarbeitet und unverzüglich gelöscht.
Es wird darauf hingewiesen, dass anonymisierte Daten und Daten, die in wissenschaftliche Auswertungen eingeflossen sind, nicht mehr gelöscht werden können.

\section*{Einwilligungserklärung}
Ich, \TextField[name=teilnehmer,width=7cm, default={Vorname Name}]{}, erkläre mich bereit,
im Rahmen der oben beschriebenen Masterarbeit an einem Interview teilzunehmen.
Ich wurde über Ziel, Ablauf und Umgang mit meinen Daten informiert.
Ich bin einverstanden, dass das Interview digital aufgezeichnet, transkribiert und für wissenschaftliche Auswertungen verwendet wird.
Weiterhin wurde ich über meine Rechte, die Freiwilligkeit der Teilnahme, die Löschfristen sowie die Möglichkeit des Widerrufs informiert.

\CheckBox[name=kontakt]{} Ich willige ein, dass meine Kontaktdaten zum Zweck möglicher Rückfragen oder einer erneuten Kontaktaufnahme im Rahmen der Studie gespeichert werden. (Bitte ankreuzen, wenn einverstanden)

\vspace{1cm}
\TextField[name=ortdatum,width=5cm]{Ort, Datum} \hspace{1cm} \TextField[name=unterschrift,width=6cm]{Unterschrift}

\end{document}
